%% =====================================================================
%% Cover letter — IEEE Access submission (standalone, PDF-ready)
%% Compila sozinho (pdfLaTeX) no Overleaf ou localmente. Não depende
%% do ieeeaccess.cls — é um documento independente.
%% =====================================================================
\documentclass[11pt]{letter}
\usepackage[T1]{fontenc}
\usepackage[utf8]{inputenc}
\usepackage[a4paper,margin=1in]{geometry}
\usepackage{parskip}

\signature{Milton Miranda Neto, on behalf of the authors\\
Universidade Federal de Uberl\^andia (UFU)\\
voidmmn@ufu.br}

\address{Milton Miranda Neto\\
Faculty of Electrical Engineering\\
Universidade Federal de Uberl\^andia (UFU)\\
Uberl\^andia, MG, Brazil\\
voidmmn@ufu.br}

\begin{document}
\begin{letter}{The Editor-in-Chief\\\textit{IEEE Access}}

\opening{Dear Editor,}

We are pleased to submit a thoroughly revised version of our manuscript,
previously ``Affective Attraction--Repulsion Value Modulation for Reactive
Inspection Navigation in Electrical Substations'' (Access-2026-33687), now
retitled ``\textbf{A Two-Layer Deliberative--Reactive Architecture with an
Affective Attraction--Repulsion Layer for Autonomous Substation Inspection},''
for consideration as a \textbf{Research Article} in \textit{IEEE Access}. A
detailed, point-by-point response to the four reviewers accompanies this
submission.

Autonomous inspection of electrical substations is typically handled either as
a fixed-coverage patrol, which classifies whatever appears without adapting its
motion to what it senses, or by purely reactive controllers, which respond only
to stimuli already within sensor range. Our work bridges these regimes with a
two-layer method: a deliberative \emph{active-sensing} patrol, built on a
standard, reproducible industrial stack (ROS~2, Nav2 with a Model Predictive
Path Integral controller, Gazebo Fortress), carries the sensors to where events
may occur, while a reactive \emph{affective} layer---named Stay Alert---makes the
robot approach anomalies to inspect them and retreat from hazards, arbitrated by
a hazard-first attraction/repulsion field.

The reviewers converged on one central point, which we have addressed
structurally rather than cosmetically. In the tabular instantiation evaluated,
the executed behavior is produced by the threshold-triggered reactive layer on
the raw affective field, while the Bellman-value term does not influence any
executed action. We have realigned the paper's title, abstract, and claims with
the mechanism that actually controls the robot, and we present the value update
as a formulation and substrate for future learned policies rather than as the
behavior driver---a concession the evidence required, which we make explicitly.

We believe the revised manuscript is a good fit for \textit{IEEE Access} for
three reasons:

\begin{enumerate}
\item \textbf{A validated, honestly bounded contribution.} Across a uniform
battery of sixteen configurations, the fixed baseline inspects no off-route
anomalies, whereas the proposed layer inspects all three in the default scenario
and under moderate noise, and consistently performs the hazard-avoidance
maneuver. We report its limits with the same care: completeness degrades under
dense clustering ($9/10$ runs complete all five events) and under high perception
noise ($4/9$), while no contact-region entries (zero collisions) were observed in
any scenario. On a hazard on the patrol route, the method contours at a
$1.71$~m minimum clearance versus the baseline's $0.71$~m, and a keepout/retreat
ablation identifies the deliberative keepout as the mechanism that resolves the
static hazard; a directed emerging-hazard trial isolates the reactive retreat's
transient safety margin, with its cost stated.

\item \textbf{Reproducibility and technical soundness.} The method runs on a
standard, openly available stack; the complete code, simulation world, launch
files, batch runner, and analysis scripts are released publicly under a tagged
release, together with the per-run data, a per-episode audit trail, a provenance
manifest, and file checksums. The evaluation uses distribution-free statistics
(Mann--Whitney~$U$, Cliff's~$\delta$, bootstrap confidence intervals, Holm
correction, and equivalence testing) appropriate to the bounded-count and
constant-variance outcomes the reviewers flagged.

\item \textbf{Transparent empirical reporting.} We state the boundary of the
value term explicitly, correct the earlier overstatements the reviewers
identified (mission time is a measured cost, not ``no significant increase'';
completeness is qualified, not universal; safety is reported as observed contacts
and measured margins, not a general guarantee), and we are candid about the
remaining limitations---chiefly that the demonstrated originality is that of
integration and analysis, since we do not isolate the affective weighting against
a heuristic supervisor using the same controller, and that validation is in
simulation.
\end{enumerate}

We confirm that this manuscript is original, has not been published previously,
and is not under consideration for publication elsewhere. All authors have
approved the submission and meet the IEEE authorship criteria. We have no
conflicts of interest to declare. In accordance with IEEE policy, the use of
AI-based tools in the preparation of this manuscript (for two conceptual figures
and for assistance with text and analysis-script editing) is disclosed in the
Acknowledgment section; the authors have reviewed and verified all such content
and take full responsibility for it.

We thank the Editor and the reviewers for an unusually careful and constructive
review, which materially improved the work, and we look forward to your decision.

\closing{Sincerely,}

\end{letter}
\end{document}
